\section{Voyager：Code-as-Action、Memory 与 Automatic Curriculum}

MineCLIP 通过 reward 把语义接入低层 policy，但长 horizon 开放探索仍需要更可组合的 action representation。Voyager 选择\term{code-as-action（代码即动作）}：LLM 不逐帧输出键盘指令，而是生成调用 Minecraft API 的 JavaScript 程序。程序可以包含循环、条件、错误处理和已有 skill 调用，从而把长动作序列压缩成可执行、可检查的符号单元。

\subsection{为什么要把视频与世界“stringify”}

互联网视频包含丰富行为，却没有干净的 action label；像素到连续控制的逆问题通常多解。课程用“stringify”表达一种接口策略：把环境状态、inventory、错误和可用 API 转成语言/代码上下文，让 LLM 在其擅长的 token space 中推理，再由工具执行并返回真实结果。

\lecturefigure{slide-26-stringify-video.jpg}{从视觉轨迹到可供语言模型处理的符号接口}{官方录像恢复幻灯片 V026；00:21:39}

\begin{knowledgebox}{读图：stringification 是接口工程，不是信息无损转换}
图中从动态视频转向文本/符号模块。转换必须选择哪些 state 暴露给模型、怎样描述空间关系、哪些 action 封装为 API。它提高可读性和可组合性，却可能丢失细粒度几何、时序与接触信息，因此更适合高层 Minecraft skill，不应被直接视为所有机器人低层控制的通用答案。
\end{knowledgebox}

\lecturefigure{slide-27-mineflayer-javascript.jpg}{Mineflayer JavaScript API 提供可执行 Minecraft 控制接口}{官方录像恢复幻灯片 V027；00:22:15}

Mineflayer 把导航、采集、制作和交互封装成程序调用。读图时先看函数参数与对象名称，再看执行结果如何返回；API 让代码可测试、可重用，也建立 safety boundary。与此同时，封装质量决定能力上限：API 不暴露的细节，LLM 无法仅靠更聪明的 prompt 补回。

\teachervoice{讲者指出视频很 rich，却没有 action label；Voyager 因此不尝试直接从每一帧恢复键盘操作，而是利用 Mineflayer 的 symbolic interface。这个选择把“理解目标”与“执行程序”连接起来，也让 syntax error、runtime error 和 task progress 都能成为模型可读反馈。}

\subsection{Voyager Data Flow 与 Iterative Prompting}

Voyager 不是一次 prompt 生成一段代码后就结束。它维护当前 state、目标、可用 skill 与执行反馈，反复让 actor 生成或修订程序；environment 则用真实状态、错误和 progress 约束模型。\term{critic（批评器）}在这里是判断任务是否完成、失败在哪里并给出改进方向的模块，不是只给一个风格化意见。

\lecturefigure{slide-28-voyager-data-flow.jpg}{Voyager 的 curriculum、code action、environment 与 memory 数据流}{官方录像恢复幻灯片 V028；00:23:12}

\begin{knowledgebox}{读图：沿闭环追踪四类状态}
从 curriculum 产生 task，actor 读取 observation 与 relevant skills，生成 code；environment 执行后返回 world state 与 error；critic 判断 progress；成功 skill 写入 library。关键不是有多少 GPT-4 调用，而是 task、execution、evaluation、memory 四类状态是否形成闭环。
\end{knowledgebox}

\lecturefigure{slide-29-code-as-action.jpg}{GPT-4 生成可执行程序控制 embodied agent}{官方录像恢复幻灯片 V029；00:23:27}

这页把传统 action token 替换为完整程序。先看 actor 输出中的函数调用、循环与检查，再看 agent 在世界中的长时间行为；代码把数百步动作压缩为一个 skill，但也引入 prompt injection、无限循环、破坏性调用和 dependency error 等软件风险，必须在 sandbox 与 API allowlist 中执行。

\lecturefigure{slide-30-iterative-prompting.jpg}{环境反馈、执行错误与自验证驱动代码迭代}{官方录像恢复幻灯片 V030；00:24:09}

\begin{warningbox}{Iterative prompting 不等于错误会自动消失}
反馈只有在可观测、可归因且及时返回时才有用。若 environment state 缺失、critic 误判、程序改变世界后无法恢复，下一轮 prompt 可能在错误前提上继续修补。系统应限制 retry budget、保留 execution trace，并把不可逆 action 单独审批。
\end{warningbox}

\begin{lstlisting}[caption={Voyager 代码生成、执行反馈与 skill memory 的概念循环}]
task = curriculum.propose(agent_state, skill_library)
while retry_budget > 0:
    skills = skill_library.retrieve(task, agent_state)
    program = actor.write_code(task, agent_state, skills)
    result = sandbox.execute(program, allowed_api=mineflayer_api)
    verdict = critic.check(task, result, agent_state)
    if verdict.success:
        skill_library.store(task, program, result.preconditions)
        break
    agent_state = result.updated_state
    actor.reflect(verdict.feedback, result.errors)
    retry_budget -= 1
\end{lstlisting}

\teachervoice{课堂把 Voyager 称为 no-gradient architecture：运行时不更新模型参数，而是通过 prompt、执行反馈和外部 memory 改善行为。这里的“learning”主要发生在 skill library 与 curriculum state 中，不能与 neural-network weight update 混为一谈。}

\subsection{Critic 与 Skill Library：把一次成功变成可复用能力}

若每个任务都从空白 prompt 开始，开放探索会重复支付相同成本。\term{skill library（技能库）}是持久化的 executable memory：保存成功程序、用途、前置条件与检索描述，使后续任务能组合旧技能。Critic 则把 environment evidence 转成完成判断和修订信号，防止仅凭代码“看起来正确”就写入 memory。

\lecturefigure{slide-31-critic-self-reflection.jpg}{Critic 根据任务进展生成 self-reflection}{官方录像恢复幻灯片 V031；00:24:45}

\begin{knowledgebox}{读图：critic 需要引用可验证状态}
有效 feedback 应说明缺少哪个物品、当前位置为何不满足条件、执行到哪一步失败，而不是笼统说“try again”。先检查 verdict 引用了哪些 environment observation，再检查建议是否能改变下一次程序。Critic 本身也可能 hallucinate，因此 success 最好由 state predicate 与语义判断共同决定。
\end{knowledgebox}

\lecturefigure{slide-32-skill-library-overview.jpg}{Skill library 持久化越来越复杂的 executable procedures}{官方录像恢复幻灯片 V032；00:25:27}

技能库的价值来自 compositionality：早期的采集、制作、导航技能可被更高层任务调用。读图时先看技能粒度是否足够稳定，再看 skill 之间的依赖；过细会造成 prompt 膨胀，过粗则难以重用。每个 skill 还需要前置条件、side effect 与 failure mode，而不只是函数名。

对 task query $q$、skill description $d_i$ 与当前状态 $x$，一个简化检索规则为：
\[
i^*=\arg\max_i\left[\operatorname{sim}(E(q,x),E(d_i))
-\lambda\,\operatorname{cost}(i,x)\right].
\]
其中 $E$ 是 embedding encoder，$\operatorname{sim}$ 衡量任务与技能语义相关性，$\operatorname{cost}(i,x)$ 表示当前状态下调用技能的资源或风险，$\lambda$ 控制相关性与成本的权衡，$i^*$ 是被检索的 skill。只用语义相似度可能取回名称相近却前置条件不满足的程序。

\lecturefigure{slide-33-retrieve-old-skills.jpg}{检索并组合旧 skill 完成更复杂的新任务}{官方录像恢复幻灯片 V033；00:26:36}

这页展示 retrieve--compose 路径：新任务先触发相关旧程序，再由 actor 生成组合逻辑。先检查旧技能是否真的被调用，再比较从零生成与复用的 token、错误和时间成本。它支持 external memory 能累积 competence，却不证明程序在环境版本变化后仍可靠；library 需要 regression test 与失效机制。

\teachervoice{讲者强调 skill library 让 Voyager 不是“一次性 ChatGPT demo”。成功代码被保存并在更难任务中检索，系统因此在不更新权重的情况下获得长期能力增长。这里最值得迁移的思想是把 agent memory 存成可执行且可验证的 artifact，而不是只存自然语言摘要。}

\subsection{Automatic Curriculum：谁来决定下一步学什么}

开放世界没有固定关卡顺序，因此还需要\term{automatic curriculum（自动课程）}：根据当前 inventory、已解锁技能、探索历史和 novelty，提出“现在可行但尚未完成”的下一任务。它同时控制难度和探索方向；若任务太难会浪费调用，太易则无法扩展能力。

\lecturefigure{slide-34-automatic-curriculum.jpg}{Automatic curriculum 沿 Minecraft 技术树扩展能力}{官方录像恢复幻灯片 V034；00:27:18}

\begin{knowledgebox}{读图：curriculum 是 task policy}
图中从基础资源走向高级物品，显示任务之间有前置依赖。先看当前 inventory 与 unlocked skills，再看候选任务是否带来新状态；好的 curriculum 不是简单追逐最高 reward，而是在 feasibility、novelty、usefulness 和 risk 之间平衡。
\end{knowledgebox}

\lecturefigure{slide-35-proposed-next-task.jpg}{Curriculum 根据状态提出可行且新颖的下一任务}{官方录像恢复幻灯片 V035；00:28:30}

任务选择可抽象为：
\[
\tau_{k+1}=\arg\max_{\tau\in\mathcal{C}(x_k)}
\bigl[\alpha F(\tau\mid x_k)+\beta N(\tau\mid M_k)+\eta U(\tau)-\rho R(\tau)\bigr].
\]
其中 $\mathcal{C}(x_k)$ 是当前状态 $x_k$ 下候选任务集合，$F$ 是 feasibility，$N$ 是相对 memory $M_k$ 的 novelty，$U$ 是未来复用价值，$R$ 是风险或预计成本，$\alpha,\beta,\eta,\rho$ 是权重。真实 Voyager 用 LLM proposal，而公式用于揭示评价维度，不表示系统显式计算了这些标量。

\subsection{闭环结果：技术树速度与地图覆盖}

将 curriculum、actor、critic 和 skill library 组装后，Voyager 才成为开放探索系统。结果页分别衡量技术树 progress 与 spatial exploration；前者接近技能/资源解锁，后者接近地图覆盖。两类指标互补，但都只是 Minecraft 内的 proxy。

\lecturefigure{slide-36-voyager-all-pieces.jpg}{Voyager 的 curriculum、prompting 与 skill library 组成闭环}{官方录像恢复幻灯片 V036；00:29:03}

\begin{importantbox}{Voyager 的最小闭环}
Curriculum 提出任务，actor 写代码，environment 返回执行证据，critic 判断进展，成功程序进入 skill library，新的 memory 又改变下一任务和下一段代码。任何一环缺失，系统都会退化为无方向探索、一次性生成、不可验证执行或无法积累的 demo。
\end{importantbox}

\lecturefigure{slide-37-open-ended-exploration-results.jpg}{Voyager 在技术树与开放探索指标上的进展曲线}{官方录像恢复幻灯片 V037；00:30:21}

\begin{knowledgebox}{读图：先看横轴预算，再看曲线含义}
横轴代表迭代或探索预算，纵轴对应解锁物品、里程碑或任务 progress；先比较 Voyager 与 baselines 在相同预算下的上升速度，再看最终平台。曲线支持 memory/curriculum/code loop 在 Minecraft 中更高效，不支持“参数模型已获得通用智能”或对真实机器人直接外推。
\end{knowledgebox}

\lecturefigure{slide-38-map-coverage.jpg}{地图覆盖比较显示开放探索的空间范围}{官方录像恢复幻灯片 V038；00:31:06}

地图中的颜色与区域表示不同 agent 的探索轨迹或覆盖范围。先比较可达面积，再检查是否只是随机游走；结合任务完成与资源获取，才能判断探索是否有用。更大覆盖支持 agent 主动访问更多状态，却不说明每个区域都被理解，也不代表学会了长期因果模型。

\teachervoice{讲者在展示结果时强调 Voyager 更快解锁 technology tree、探索更大地图；同时这仍是 Minecraft progress。把指标放回环境边界很重要：地图覆盖和物品数量是开放探索 proxy，不是 general intelligence 的完整测量。}

\subsection{本章小结}

Voyager 用 code-as-action 把长行为压缩为可执行程序，以 environment feedback 和 critic 反复修订，用 skill library 积累 executable memory，再用 automatic curriculum 决定下一任务。它展示了一种无需更新模型权重的持续改进架构，但能力仍受 API、critic、memory hygiene、task proposal 与 Minecraft 指标约束。

